\documentclass[a4paper,11pt]{article}
\usepackage[T2A]{fontenc}
\usepackage[utf8]{inputenc}
\usepackage[russian]{babel}
\usepackage{amsmath,amssymb,mathtools}
\renewcommand{\familydefault}{\sfdefault}
\usepackage{icomma}
\usepackage[margin=23mm,top=22mm,bottom=22mm,headheight=22pt]{geometry}
\usepackage{microtype}
\usepackage{tikz}
\usetikzlibrary{arrows.meta,positioning,shapes.geometric}
\usepackage{tabularx,booktabs,array}
\usepackage{enumitem}
\usepackage{fancyhdr}
\usepackage{xcolor}
\usepackage{hyperref}
\usepackage{parskip}
\definecolor{accent}{HTML}{008080}
\definecolor{darkaccent}{HTML}{005656}
\definecolor{textgray}{HTML}{333333}
\definecolor{muted}{HTML}{6D7476}
\hypersetup{hidelinks,pdftitle={Чек-лист. Логарифмы и их преобразования},pdfauthor={Лёвин Артём Александрович}}
\setlength{\parindent}{0pt}
\setlength{\parskip}{2.5pt}
\linespread{1.035}
\setlist[enumerate]{leftmargin=*,itemsep=5pt,topsep=5pt}
\setlist[itemize]{leftmargin=*,itemsep=2.5pt,topsep=3pt}
\pagestyle{fancy}
\fancyhf{}
\fancyhead[L]{\footnotesize\color{muted} Логарифмы и их преобразования}
\fancyhead[R]{\footnotesize\color{muted} Даниил Хрястов}
\fancyfoot[C]{\footnotesize\color{muted}\thepage}
\renewcommand{\headrulewidth}{0.35pt}
\renewcommand{\headrule}{\hbox to\headwidth{\color{accent}\leaders\hrule height \headrulewidth\hfill}}
\newcommand{\sect}[2]{\vspace{8pt}{\large\color{darkaccent}\bfseries #1.\ #2}\par\vspace{2pt}\color{accent}\rule{\linewidth}{0.45pt}\color{textgray}\vspace{2pt}}
\newcommand{\ex}[1]{\smallskip{\color{darkaccent}\bfseries Пример #1.}\enspace}
\newcommand{\idea}[1]{{\color{darkaccent}\bfseries Ключевая идея.}\enspace #1}
\newcommand{\ans}[1]{\(\displaystyle #1\)}
\begin{document}
\color{textgray}
\thispagestyle{empty}
\begin{center}
\vspace*{29mm}
{\normalsize\color{darkaccent} МАТЕМАТИКА}\par
\vspace{15mm}
{\LARGE Чек-лист}\par
\vspace{12mm}
{\Large\color{accent} Логарифмы: свойства,}\par
{\Large\color{accent} преобразования и вычисления}\par
\vspace{15mm}
{\normalsize Подготовка к ЕГЭ по профильной математике}\par
\vfill
{\normalsize Даниил Хрястов}\par
\vspace{9mm}
{\normalsize 11 ноября 2026 г.}\par
\vspace{18mm}
{\small Лёвин Артём Александрович}\par
{\small эксклюзивно для Даниила Хрястова}\par
\vspace{17mm}
\begin{tikzpicture}
\draw[accent,line width=.95pt] (-3.4,0) .. controls (-2.0,.48) and (-1.0,-.38) .. (0,0)
.. controls (1.1,.46) and (2.1,-.40) .. (3.4,0);
\fill[darkaccent] (0,0) circle (1pt);
\end{tikzpicture}
\end{center}
\clearpage

\sect{1}{Главная идея и связь со степенями}
Логарифм отвечает на вопрос: \emph{«В какую степень нужно возвести основание, чтобы получить заданное число?»}
\[
\boxed{\log_a b=c\quad\Longleftrightarrow\quad a^c=b},
\qquad a>0,\quad a\ne1,\quad b>0.
\]
\idea{При вычислении логарифмов дроби, корни и составные числа удобно сначала представить степенями. Формулы полезно применять в обоих направлениях.}

\ex{1} Вычислите \(\log_2 8\):
\[
\log_2 8=\log_2 2^3=3.
\]
\ex{2} Вычислите \(\displaystyle\log_2\frac1{\sqrt[3]{8}}\):
\[
\frac1{\sqrt[3]{8}}=\frac1{\sqrt[3]{2^3}}=2^{-1},
\qquad\log_2 2^{-1}=-1.
\]

\textbf{Повторение предыдущей темы.} Если график \(f(x)=a^x+b\) проходит через \((0;-2)\) и \((1;-1)\), то
\[
1+b=-2,\qquad a+b=-1\quad\Longrightarrow\quad b=-3,\ a=2,
\qquad f(4)=2^4-3=13.
\]
Знание свойств степеней непосредственно понадобится и в логарифмах.

\sect{2}{Обязательная проверка допустимости}
Перед преобразованиями убедитесь, что \textbf{основание положительно и не равно единице}, а \textbf{аргумент строго положителен}:
\[
\log_{g(x)} h(x):\quad g(x)>0,\quad g(x)\ne1,\quad h(x)>0.
\]
\ex{3} Выражение \(\log_2(x-5)\) определено при \(x>5\).
Выражение \(\log_2\bigl((x-5)^2\bigr)\) определено при \(x\ne5\): квадрат положителен для всех ненулевых \(x-5\).

\sect{3}{Степень аргумента: вынесение показателя}
Если \(a>0,\ a\ne1,\ b>0\), то для любого действительного \(r\):
\[
\boxed{\log_a(b^r)=r\log_a b}
\quad\Longleftrightarrow\quad
r\log_a b=\log_a(b^r).
\]
\ex{4} \(\log_3 81=\log_3 3^4=4\); а \(2\log_5 3=\log_5 9\).

\textbf{Особый случай чётной целой степени.} При \(b\ne0\) и целом \(n\ne0\):
\[
\boxed{\log_a(b^{2n})=2n\log_a|b|.}
\]
Модуль обязателен, поскольку исходное \(b\) может быть отрицательным.
\ex{5} При \(x\ne5\):
\[
\log_2\bigl((x-5)^2\bigr)=2\log_2|x-5|.
\]
Без модуля правая часть при \(x<5\) оказалась бы не определена, хотя левая часть существует.

\clearpage
\sect{4}{Степень основания: обратный показатель}
При \(a>0,\ a\ne1,\ b>0,\ r\ne0\):
\[
\boxed{\log_{a^r}b=\frac1r\log_a b.}
\]
\textbf{Различайте операции:} степень \emph{аргумента} выносится множителем \(r\); степень \emph{основания} даёт обратный множитель \(1/r\).

\ex{6} Сначала приводим основание к степени:
\[
\log_{2^{3/5}}2=\frac53\log_2 2=\frac53,
\qquad
\log_{\sqrt[3]2}2=3.
\]
\ex{7} Если основание содержит четвёртую степень и может быть отрицательная внутренняя величина, то
\[
\log_{(x-3)^4}7=\frac14\log_{|x-3|}7,
\quad x\notin\{2,3,4\}.
\]
Здесь \(x=3\) обращает основание в нуль, а \(x=2,4\) делают основание равным единице.

\sect{5}{Сложение и вычитание логарифмов}
При одинаковом допустимом основании и положительных \(b,c\):
\[
\boxed{\log_a b+\log_a c=\log_a(bc)},\qquad
\boxed{\log_a b-\log_a c=\log_a\frac bc}.
\]
Обе формулы работают в обратную сторону. Если перед логарифмом есть множитель, сначала его можно занести в степень аргумента. Для объединения основания должны совпадать.

\ex{8} Объединим три логарифма:
\[
\log_2 5+\log_2 7-\log_2 3
=\log_2\frac{5\cdot7}{3}=\log_2\frac{35}{3}.
\]
\ex{9} Превратим число в логарифм нужного основания:
\[
3=3\log_2 2=\log_2 2^3=\log_2 8,
\]
\[
3+\log_2 7=\log_2 8+\log_2 7=\log_2 56.
\]
Вообще, при допустимом \(a\): \(t=\log_a a^t\). Это полезно для целых и дробных коэффициентов и даже для переменной \(t\).

\ex{10} Если основания различаются, приведём их к одному:
\[
\log_2 3+\log_4 5
=\log_2 3+\frac12\log_2 5
=\log_2(3\sqrt5).
\]

\sect{6}{Основное логарифмическое тождество}
\[
\boxed{a^{\log_a b}=b},\qquad a>0,\ a\ne1,\ b>0.
\]
\ex{11} Лишний коэффициент перед логарифмом переносим в его аргумент:
\[
2^{3\log_2 5}=2^{\log_2 5^3}=5^3=125.
\]
\ex{12} При сложении показателей нужно использовать свойства степеней:
\[
2^{3+\log_2 5}=2^3\cdot 2^{\log_2 5}=8\cdot5=40.
\]
\idea{В степенном выражении ищите совпадение основания степени с основанием логарифма. Если их нет, попробуйте привести оба числа к общему основанию.}

\sect{7}{Переход к новому основанию и обратные логарифмы}
Если \(a,b,c>0\), \(a\ne1\), \(c\ne1\), то
\[
\boxed{\log_a b=\frac{\log_c b}{\log_c a}}.
\]
Частный случай — обмен основания и аргумента:
\[
\boxed{\log_a b=\frac{1}{\log_b a}},\qquad b\ne1.
\]
\ex{13} \(\log_2 3\cdot\log_3 2=1\).
\ex{14} В выражении \(4^{\log_8 5}\) обе основы приводим к двойке:
\[
4^{\log_8 5}=\left(2^2\right)^{\frac13\log_2 5}
=2^{\frac23\log_2 5}=5^{2/3}.
\]
\textbf{Полезная симметрия:} при \(a,c>0\) и \(b>0,\ b\ne1\):
\[
\boxed{a^{\log_b c}=c^{\log_b a}}.
\]
Она особенно удобна, когда числа в основании и аргументе хочется поменять местами.

\sect{8}{Составные выражения: выбирайте подходящий маршрут}
\ex{15} Здесь ключевую работу выполняют свойства степеней:
\[
\frac{9^{\log_5 50}}{9^{\log_5 2}}
=9^{\log_5 50-\log_5 2}
=9^{\log_5 25}=9^2=81.
\]
Другой путь — применить симметрию к числителю и знаменателю, а затем разделить степени с одинаковым показателем.

\ex{16} Вложенный логарифм вычисляйте \emph{изнутри наружу}:
\[
\log_2\bigl(\log_3 81\bigr)=\log_2 4=2.
\]
\ex{17} Квадрат \emph{всего логарифма} — отдельная операция:
\[
\left(\log_{\sqrt7}49\right)^2
=\left(\frac{\log_7 7^2}{\log_7 7^{1/2}}\right)^2
=4^2=16.
\]
Нельзя переносить внешний квадрат на аргумент по формуле вынесения степени.

\clearpage
\sect{9}{Частые ошибки и самопроверка}
\begin{itemize}
\item \(\log_a(b+c)\) нельзя превращать в \(\log_a b+\log_a c\). Формулы суммы \emph{аргументов} нет.
\item При переходе от \(\log_{a^r}b\) появляется множитель \(1/r\), а не \(r\).
\item При выносе \emph{чётной} степени выражения с возможным отрицательным значением сохраняйте модуль.
\item \(\left(\log_a b\right)^2\) и \(\log_a(b^2)\) обозначают разные выражения.
\item При сложении/вычитании объединяйте только логарифмы с одинаковым основанием и учитывайте множители перед ними.
\item Всегда проверяйте ОДЗ исходного выражения: преобразование не должно расширять её незаметно.
\end{itemize}


\vspace{7mm}
{\color{darkaccent}\bfseries Итоговая памятка}\par
\smallskip
\begin{tabularx}{\linewidth}{@{}p{.39\linewidth}X@{}}
\toprule
Сигнал в задаче & Полезное действие\\
\midrule
Корень, дробь, составное число & Представьте в степенном виде.\\
Степень аргумента & Вынесите показатель; проверьте модуль.\\
Степень основания & Вынесите обратный показатель.\\
Сумма/разность логарифмов & Приведите основания к одному и объедините аргументы.\\
Основание степени совпадает с основанием логарифма & Примените \(a^{\log_a b}=b\).\\
Основания мешают друг другу & Перейдите к новому основанию.\\
Логарифм внутри логарифма & Вычисляйте изнутри наружу.\\
\bottomrule
\end{tabularx}

\clearpage
\begin{center}
{\Large\color{darkaccent}\bfseries Алгоритм: как преобразовывать логарифм}
\end{center}
\vspace{11mm}
\begin{center}
\begin{tikzpicture}[
 >=Latex,
 node distance=9mm and 18mm,
 action/.style={rectangle,rounded corners=2pt,draw=darkaccent,align=center,text width=9.0cm,minimum height=1.14cm,font=\small,inner sep=4pt},
 decision/.style={diamond,aspect=3.5,draw=accent,align=center,text width=5.7cm,minimum height=1.45cm,font=\small,inner sep=2pt},
 endpoint/.style={ellipse,draw=darkaccent,align=center,text width=5.6cm,minimum height=.85cm,font=\small},
 link/.style={-{Latex[length=2.8mm]},thick,draw=darkaccent}
]
\node[endpoint] (start) {Получено выражение с логарифмами};
\node[action,below=of start] (domain) {Проверьте ОДЗ: основание $>0$, $\ne1$; аргумент $>0$};
\node[decision,below=of domain] (powers) {Есть дроби, корни или составные числа?};
\node[action,below=of powers] (normalize) {Приведите числа к степеням; выделите одинаковые основания};
\node[action,below=of normalize] (formulas) {Примените формулы: показатели, сложение, вычитание, переход к основанию};
\node[decision,below=of formulas] (done) {Можно вычислить или упростить?};
\node[action,below=of done] (verify) {Проверьте равносильность, модуль и ОДЗ};
\node[endpoint,below=of verify] (end) {Запишите результат};
\node[action,right=14mm of done,text width=3.9cm] (retry) {Смените стратегию: степенное/логарифмическое представление};
\draw[link] (start)--(domain);
\draw[link] (domain)--(powers);
\draw[link] (powers)--node[right,font=\footnotesize]{да}(normalize);
\draw[link] (normalize)--(formulas);
\draw[link] (formulas)--(done);
\draw[link] (done)--node[right,font=\footnotesize]{да}(verify);
\draw[link] (verify)--(end);
\draw[link] (done.east)--node[above,font=\footnotesize]{нет}(retry.west);
\draw[link] (retry.north) |- (formulas.east);
\draw[link] (powers.west) -- ++(-1.05,0) |- (formulas.west) node[pos=.13,left,font=\footnotesize]{нет};
\end{tikzpicture}
\end{center}
\clearpage

\sect{10}{Тренировочный блок · Путь А}
Решите \textbf{10 заданий} по порядку. Записывайте промежуточные преобразования; у выражений с переменной указывайте ограничения.

{\color{darkaccent}\bfseries Базовый уровень}\vspace{-1mm}
\begin{enumerate}[label=\arabic*.,itemsep=7pt]
\item Вычислите: \(\log_2 64\).
\item Вычислите: \(\displaystyle\log_5\frac1{\sqrt{25}}\).
\item Вычислите: \(\displaystyle\log_{3^{2/3}}9\).
\item Преобразуйте с использованием модуля и укажите ОДЗ: \(\log_7\bigl((x-3)^2\bigr)\).
\end{enumerate}
{\color{darkaccent}\bfseries Уровень выше среднего}\vspace{-1mm}
\begin{enumerate}[label=\arabic*.,start=5,itemsep=8pt]
\item Представьте одним логарифмом: \(3-\log_2 5\).
\item Представьте одним логарифмом с основанием \(2\): \(\log_4 7+\log_2 3\).
\item Вычислите: \(\log_5 7\cdot\log_7 125\).
\item Вычислите: \(8^{\log_4 9}\).
\end{enumerate}
{\color{darkaccent}\bfseries Повышенный уровень}\vspace{-1mm}
\begin{enumerate}[label=\arabic*.,start=9,itemsep=8pt]
\item Вычислите: \(\displaystyle\frac{11^{\log_5 75}}{11^{\log_5 3}}\).
\item Вычислите, \textbf{не раскрывая скобки}: \(\bigl(1-\log_3 15\bigr)\bigl(1-\log_5 15\bigr)\).
\end{enumerate}
\vfill
\small{\color{muted} Подсказка к самопроверке: для заданий 4--6 важно сохранить ОДЗ и правильную структуру логарифма; для заданий 7--10 полезны обратные преобразования и свойства степеней.}
\clearpage

\sect{11}{Ответы к тренировочному блоку}
\renewcommand{\arraystretch}{1.65}
\begin{center}
\begin{tabularx}{\linewidth}{@{}>{\bfseries}p{10mm}X@{}}
\toprule
№ & Ответ / корректная форма преобразования\\
\midrule
1 & \(6\)\\
2 & \(-1\)\\
3 & \(3\)\\
4 & \(2\log_7|x-3|\), при \(x\ne3\)\\
5 & \(\log_2\dfrac85\)\\
6 & \(\log_2(3\sqrt7)\)\\
7 & \(3\)\\
8 & \(27\)\\
9 & \(121\)\\
10 & \(1\)\\
\bottomrule
\end{tabularx}
\end{center}
\vfill
\begin{center}
\footnotesize\color{muted} Лёвин Артём Александрович эксклюзивно для Даниила Хрястова
\end{center}
\end{document}
